% Section content template for individual Educative lessons
% This file contains the actual content for a specific section
% Generated from Educative API section content

% Section: Design of a Key-Value Store
% Section ID: 5610927277211648
% Section Slug: design-of-a-key-value-store
% Generated: 2025-12-06T22:41:32.490670

\subsection{Requirements}\label{MAsBeqDWYmLPkkIRvY6PP}
\\
Let's list the requirements of designing a key-value store to overcome the problems of traditional databases.

\subsubsection{Functional requirements}\label{CtIw12xjK_DrJYMPBBwg7}
\\
Typically, key-value stores are expected to offer functions such as \texttt{get} and \texttt{put}. However, what sets this particular key-value store system apart is its distinct characteristics, explained as follows:

\begin{itemize}
\item
\protect\phantomsection\label{pRfuKVAGSDUZlL7waSmnY}
\textbf{Configurable service}: Some applications might have a tendency to trade strong consistency for higher availability. We need to provide a configurable service so that different applications could use a range of consistency models. We need tight control over the trade-offs between availability, consistency, cost-effectiveness, and performance.
\end{itemize}
\begin{quote}
\textbf{Note:} Such configurations can only be performed when instantiating a new key-value store instance and cannot be changed dynamically when the system is operational.
\end{quote}
\begin{itemize}
\item
\protect\phantomsection\label{dtEJxr50s6C2K6ripTP-M}
\textbf{Ability to always write (when we picked ``A'' over ``C'' in the context of CAP)}: The applications should always have the ability to write into the key-value storage. If the user wants strong consistency, this requirement might not always be fulfilled due to the implications of the CAP theorem.
\end{itemize}
\begin{quote}
\textbf{Note:} The context of the problem determines what will be classified as a functional requirement and what will be classified as non-functional. For example, the ability to always write (high availability) is a functional requirement for Amazon's shopping cart application, while in other cases, high availability may be considered a non-functional requirement. Drawing inspiration from Amazon's Dynamo key-value store, we can categorize the ability to always write as a functional requirement.
\end{quote}
\begin{itemize}
\item
\protect\phantomsection\label{r1EVP-bjuo1uq43aQZQtD}
\textbf{Hardware heterogeneity}: We want to add new servers with different and higher capacities, seamlessly, to our cluster without changing or upgrading existing servers. Our system should be able to accommodate and leverage different capacity servers, ensuring correct core functionality (\texttt{get} and \texttt{put} data) while balancing the workload distribution according to each server's capacity. This calls for a peer-to-peer design with no distinguished nodes.
\end{itemize}

\subsubsection{Non-functional requirements}\label{syLjws3-U6lL2Ggq7SfMG}
\\
The non-functional requirements are as follows:

\begin{itemize}
\item
\protect\phantomsection\label{C7pTNDPql04l_7HeBw6l9}
\textbf{Scalability}: Key-value stores should run on tens of thousands of servers distributed across the globe. Incremental scalability is highly desirable. We should add or remove the servers as needed with minimal to no disruption to the service availability. Moreover, our system should be able to handle an enormous number of users of the key-value store.
\item
\protect\phantomsection\label{J6cNFbQ09znjW-WWFGeXk}
\textbf{Fault tolerance}: The key-value store should operate uninterrupted despite failures in servers or their components.
\end{itemize}
\\
Now that you understand the functional and non-functional requirements of a key-value store, what do you think are the key differences between key-value stores and traditional databases?
\\
In what scenarios are key-value stores particularly advantageous?

\textbf{Question:}

Differences between key-value stores and traditional databases

\textbf{Reference Answer:}

Differences: Key-value stores focus on simple retrieval by key, don't have complex query languages, and trade strict consistency for availability and scalability in many cases. Key -value stores have unstructured data. Advantages
: Key-value stores excel at handling massive amounts of data, high-speed data lookups, and situations where the data structure may be flexible or not fully known in advance.

Running a key-value store on a single server limits scalability and reliability. No single machine can handle large-scale data and query demands.
\\
If that server fails, the entire system goes down. Distributing data across multiple servers helps scale storage and ensures continuous availability.

\subsection{Assumptions}\label{h5RCxPoNnAWysIfoPtcOL}
\\
We'll assume the following to keep our design simple:

\begin{itemize}
\item
\protect\phantomsection\label{z0xL0yqAuFT3IPtof776g}
The data centers hosting the service are trusted (non-hostile).
\item
\protect\phantomsection\label{vrPgmWkqrYo0vnMMjYB1y}
All the required authentication and authorization are already completed.
\item
\protect\phantomsection\label{I5_3M9uuEU4pPx1fV7p_V}
User requests and responses are relayed over HTTPS.
\end{itemize}

\subsection{API design}\label{OLL9LcibwHUkH5Nx63fS3}
\\
Key-value stores, like ordinary hash tables, provide two primary functions, which are \texttt{get} and \texttt{put}.
\\
Let's look at the API design.
\\
\textbf{The} \textbf{\texttt{get}} \textbf{function}
\\
The API call to get a value should look like this:

\begin{verbatim}
	get(key)
\end{verbatim}


We return the associated value on the basis of the parameter \texttt{key}. When data is replicated, it locates the object replica associated with a specific key that's hidden from the end user. It 's done by the system if the store is configured with a weaker data consistency model. For example, in eventual consistency, there might be more than one value returned against a key.

\begin{table}[htbp]
\centering
\small
\adjustbox{max width=\textwidth}{
\begin{tabular}{|p{0.303\textwidth}|p{0.647\textwidth}|}
\hline
\textbf{Parameter} & \textbf{Description} \\
\hline
\hline
\texttt{key} & It's the \texttt{key} against which we want to get \texttt{value}. \\
\hline
\end{tabular}
}
\end{table}

\textbf{The} \textbf{\texttt{put}} \textbf{function}
\\
The API call to put the value into the system should look like this:

\begin{verbatim}
	put(key, value)
\end{verbatim}

It stores the \texttt{value} associated with the \texttt{key}. The system automatically determines where data should be placed. Additionally , the system often keeps metadata about the stored object. Such metadata can include the version of the object.

\begin{table}[htbp]
\centering
\small
\adjustbox{max width=\textwidth}{
\begin{tabular}{|p{0.303\textwidth}|p{0.647\textwidth}|}
\hline
\textbf{Parameter} & \textbf{Description} \\
\hline
\hline
\texttt{key} & It\textquotesingle s the \texttt{key} against which we have to store \texttt{value}. \\
\hline
\texttt{value} & It\textquotesingle s the object to be stored against the \texttt{key}. \\
\hline
\end{tabular}
}
\end{table}
\\
\\
When performing data integrity checks, we often store hashes of values (and sometimes values along with their associated key) as metadata. Should these hashes be generated before or after data compression or encryption, and why?

\textbf{Question:}

Best practice for data integrity checks

\textbf{Reference Answer:}

In most data integrity use cases, hashes should be generated after compression but before encryption.

\begin{itemize}
\tightlist
\item
\textbf{Hashing before encryption} is recommended because encryption produces ciphertext, which may vary across operations due to salting or IVs, even if the underlying data is the same. Hashing encrypted data would make integrity checks unreliable.
\item
\textbf{Hashing after compression} is preferred because compression reduces data size without changing meaning, and it avoids hashing redundant data unnecessarily. So, the typical and most effective sequence is: Compress → Hash → Encrypt. This ensures integrity is validated on the actual content, not the encrypted form, and it supports consistent verification during both storage and retrieval operations. \textbf{Example scenario:} Suppose you store daily database backups in a cloud archive. First, you compress each backup to reduce its size. Then, you generate a hash of the compressed file and store it as metadata. Finally, you encrypt the file before uploading it. Later, when retrieving the backup, you decrypt it and re-hash the compressed file to compare with the stored hash. This way, you verify that the actual content is intact, regardless of encryption changes or storage format.
\end{itemize}

\subsubsection{Data type}\label{wsn66H_rub8h6Nq3PHCLX}
\\
The key is often a primary key in a key-value store, while the value can be any arbitrary binary data.
\begin{quote}
\textbf{Note:} Dynamo uses MD5 hashes on the key to generate a 128-bit identifier. These identifiers help the system determine which server node will be responsible for this specific key.
\end{quote}
\\
In the next lesson, we'll learn how to design our key-value store. First , we'll focus on adding scalability, replication, and versioning of our data to our system. Then
, we'll ensure the functional requirements and make our system fault-tolerant. We 'll fulfill a few of our non-functional requirements first because implementing our functional requirements depends on the method chosen for scalability.
\begin{quote}
\textbf{Note:} This chapter is based on Dynamo(Amazon's Highly Available Key-value), which is an influential work in the domain of key-value stores.
\end{quote}

% End of section content